% =============================================================================
% Chapter 3: System Design  (10–12 pages target)
% =============================================================================
\chapter{System Design}
\label{chap:design}

\section*{Introduction}

This chapter presents the architectural and object-oriented design of \eduspace{}.
It begins by justifying the choice of a microservices architectural paradigm and describing the overall system architecture, including the service decomposition, communication patterns, and infrastructure components.
The chapter then covers the object-oriented design artifacts: class diagrams, the persistence model, the attribute dictionary, and interaction modeling.

% =============================================================================
\section{General Architecture}
\label{sec:architecture}

% ─────────────────────────────────────────────────────────────────────────────
\subsection{Adopted Architectural Paradigm}
\label{sec:paradigm}

A fundamental design decision was whether to adopt a monolithic or a microservices architecture.
\tabref{tab:mono-vs-micro} summarizes the trade-offs between the two paradigms.

\begin{table}[H]
  \centering
  \caption{Comparison of monolithic and microservices architectures.}
  \label{tab:mono-vs-micro}
  \small
  \begin{tabularx}{\textwidth}{@{} l X X @{}}
    \toprule
    \textbf{Criterion} & \textbf{Monolithic} & \textbf{Microservices} \\
    \midrule
    Deployment        & Single unit; any change requires redeploying the entire application & Each service is deployed independently \\
    Database          & Single shared database; tight coupling between modules & Database-per-service; no shared state \\
    Scalability       & Entire application must be scaled uniformly & Individual services can be scaled based on load \\
    Fault isolation   & A failure in one module can bring down the entire system & Failures are contained within the affected service \\
    Complexity        & Lower operational complexity & Higher operational complexity (networking, orchestration) \\
    \bottomrule
  \end{tabularx}
\end{table}

The microservices approach was selected for \eduspace{} because the platform's functionality naturally decomposes into distinct bounded contexts (user management, classrooms, content, chat, notifications, search), each with its own data model and scaling characteristics~\cite{newman2021microservices}.
The database-per-service pattern~\cite{richardson2018microservices} prevents tight coupling, and different services have fundamentally different runtime needs---the communication service requires persistent WebSocket connections, the search service relies on Elasticsearch, and the file service is entirely stateless.

The backend is decomposed into seven domain services plus an API Gateway.
\tabref{tab:services} lists each service with its responsibility.

\begin{table}[H]
  \centering
  \caption{Service decomposition of the \eduspace{} backend.}
  \label{tab:services}
  \small
  \begin{tabularx}{\textwidth}{@{} l c c X @{}}
    \toprule
    \textbf{Service} & \textbf{Port} & \textbf{Database} & \textbf{Responsibility} \\
    \midrule
    API Gateway            & 3001 & ---               & JWT verification, rate limiting, request routing \\
    User Service           & 3002 & users\_db         & Authentication, profiles, password reset \\
    Class Service          & 3003 & classes\_db       & Classrooms, members, roles, chapters \\
    Content Service        & 3004 & content\_db       & Posts, comments, submissions, grading \\
    Communication Service  & 3005 & communication\_db & Real-time group chat, WebSocket, messages \\
    Notification Service   & 3006 & notifications\_db & Event-driven notifications, email delivery \\
    Search Service         & 3007 & --- (Elasticsearch) & Full-text search, post indexing \\
    File Service           & 3010 & --- (stateless)   & File upload/download via MinIO \\
    \bottomrule
  \end{tabularx}
\end{table}

Services communicate through two complementary patterns.
\textbf{Synchronous REST calls} are used when an immediate response is required---for example, the Content Service calls the Class Service to verify membership before any post operation.
\textbf{Asynchronous events via RabbitMQ} are used for fire-and-forget notifications: when a post is created, the Content Service publishes a \texttt{post.created} event consumed by the Notification and Search services, decoupling them temporally and preventing cascading failures.

All HTTP requests follow the path: React SPA $\rightarrow$ NGINX (port 5000) $\rightarrow$ API Gateway (JWT verification, rate limiting) $\rightarrow$ downstream service.
WebSocket connections for chat bypass the gateway and are proxied by NGINX directly to the Communication Service.
The infrastructure layer, orchestrated by Docker Compose, includes PostgreSQL (five databases), Redis, RabbitMQ, MinIO, Elasticsearch, and NGINX~\cite{postgresql2024docs,redis2024docs,rabbitmq2024docs,minio2024docs,elasticsearch2024docs,nginx2024docs}.

% ─────────────────────────────────────────────────────────────────────────────
\subsection{Functional Architecture (Site Map)}
\label{sec:sitemap}

\figref{fig:sitemap} presents the site map of \eduspace{}, showing the navigation structure between public pages (landing, login, registration, password reset) and protected pages (dashboard, calendar, to-do, settings, and classroom views).

\begin{figure}[H]
  \centering
  \includegraphics[width=\textwidth,height=0.85\textheight,keepaspectratio]{diagrams/sitemap.pdf}
  \caption{Site map of the \eduspace{} platform.}
  \label{fig:sitemap}
\end{figure}

% =============================================================================
\section{Object-Oriented Design}
\label{sec:oo-design}

% ─────────────────────────────────────────────────────────────────────────────
\subsection{Design Class Diagrams}
\label{sec:class-diagrams}

\figref{fig:class-design} presents the design class diagram, organized by service boundary.
Each service exposes controllers and data transfer objects (DTOs).
Inter-service dependencies are shown as REST calls (e.g., the Content Service calls the Class Service to verify membership).

\begin{figure}[H]
  \centering
  \rotatebox{90}{\includegraphics[width=0.85\textheight,height=\textwidth,keepaspectratio]{diagrams/class-design.pdf}}
  \caption{Design class diagram of \eduspace{}, organized by microservice.}
  \label{fig:class-design}
\end{figure}

% ─────────────────────────────────────────────────────────────────────────────
\subsection{Persistence Class Diagram}
\label{sec:persistence-diagram}

\figref{fig:class-persistence} presents the persistence class diagram, showing all entities across the five databases.
Each database boundary is color-coded: \texttt{users\_db}, \texttt{classes\_db}, \texttt{content\_db}, \texttt{communication\_db}, and \texttt{notifications\_db}.
Cross-database references (e.g., \texttt{class\_id} in the Content Service) are maintained as logical foreign keys, since each service owns its own database.

\begin{figure}[H]
  \centering
  \rotatebox{90}{\includegraphics[width=0.85\textheight,height=\textwidth,keepaspectratio]{diagrams/class-persistence.pdf}}
  \caption{Persistence class diagram showing all entities across the five databases.}
  \label{fig:class-persistence}
\end{figure}

% ─────────────────────────────────────────────────────────────────────────────
\subsection{Attribute Dictionary}
\label{sec:attribute-dictionary}

The following tables describe the key attributes of each entity, derived from the Prisma schema definitions.

\begin{longtable}{@{} l l l p{5.5cm} @{}}
  \caption{Attribute dictionary --- \texttt{users\_db}.} \label{tab:dict-users} \\
  \toprule
  \textbf{Entity} & \textbf{Attribute} & \textbf{Type} & \textbf{Description} \\
  \midrule
  \endfirsthead
  \toprule
  \textbf{Entity} & \textbf{Attribute} & \textbf{Type} & \textbf{Description} \\
  \midrule
  \endhead
  User & userId & UUID (PK) & Unique identifier \\
       & email & VARCHAR (UNIQUE) & User email address \\
       & password & VARCHAR & Bcrypt-hashed password \\
       & isVerified & BOOLEAN & Email verification status \\
       & verificationCode & VARCHAR? & 6-digit email verification code \\
       & codeExpiresAt & TIMESTAMP? & Expiry of verification code \\
  \midrule
  User\_profile & user\_id & UUID (PK, FK) & References User \\
               & first\_name & VARCHAR & First name \\
               & last\_name & VARCHAR & Last name \\
               & avatar\_url & VARCHAR? & URL to avatar in MinIO \\
               & phone\_number & VARCHAR? & Optional phone number \\
               & timezone & VARCHAR? & User timezone preference \\
  \midrule
  Password\_reset & id & UUID (PK) & Unique identifier \\
                 & user\_id & UUID (FK) & References User \\
                 & reset\_token\_hash & VARCHAR (UNIQUE) & SHA-256 hash of reset token \\
                 & expires\_at & TIMESTAMP & Token expiry time \\
                 & used\_at & TIMESTAMP? & When the token was used \\
  \bottomrule
\end{longtable}

\begin{longtable}{@{} l l l p{5.5cm} @{}}
  \caption{Attribute dictionary --- \texttt{classes\_db}.} \label{tab:dict-classes} \\
  \toprule
  \textbf{Entity} & \textbf{Attribute} & \textbf{Type} & \textbf{Description} \\
  \midrule
  \endfirsthead
  \toprule
  \textbf{Entity} & \textbf{Attribute} & \textbf{Type} & \textbf{Description} \\
  \midrule
  \endhead
  Classroom & id & UUID (PK) & Unique identifier \\
            & name & VARCHAR & Classroom display name \\
            & type & ENUM & TEACHING or FRIENDLY \\
            & class\_code & VARCHAR(6) (UNIQUE) & 6-char join code \\
            & description & VARCHAR? & Optional description \\
            & subject & VARCHAR? & Subject or course name \\
            & section & VARCHAR? & Section/group label \\
            & cover\_image & VARCHAR? & Cover image file key \\
            & chat\_enabled & BOOLEAN & Whether chat is active \\
            & creator\_id & UUID & User who created the classroom \\
  \midrule
  Member & id & UUID (PK) & Unique identifier \\
         & class\_id & UUID (FK) & References Classroom \\
         & user\_id & UUID & References User (logical FK) \\
         & role & ENUM & ADMIN or MEMBER \\
         & joined\_at & TIMESTAMP & When the user joined \\
  \midrule
  Chapter & id & UUID (PK) & Unique identifier \\
          & class\_id & UUID (FK) & References Classroom \\
          & name & VARCHAR & Chapter title \\
          & position & INT & Display order (0-based) \\
  \bottomrule
\end{longtable}

\begin{longtable}{@{} l l l p{5.5cm} @{}}
  \caption{Attribute dictionary --- \texttt{content\_db}.} \label{tab:dict-content} \\
  \toprule
  \textbf{Entity} & \textbf{Attribute} & \textbf{Type} & \textbf{Description} \\
  \midrule
  \endfirsthead
  \toprule
  \textbf{Entity} & \textbf{Attribute} & \textbf{Type} & \textbf{Description} \\
  \midrule
  \endhead
  Post & id & UUID (PK) & Unique identifier \\
       & class\_id & UUID & Classroom (logical FK) \\
       & chapter\_id & UUID & Chapter (logical FK) \\
       & author\_id & UUID & Author (logical FK) \\
       & title & VARCHAR & Post title \\
       & content & TEXT? & Optional body text \\
       & type & ENUM & Post type (5 values) \\
       & study\_material\_type & ENUM? & Sub-type for STUDY\_MATERIAL \\
       & due\_date & TIMESTAMP? & Assignment due date \\
       & max\_points & INT? & Maximum grade points \\
  \midrule
  Attachment & id & UUID (PK) & Unique identifier \\
             & post\_id & UUID (FK) & References Post \\
             & file\_key & VARCHAR & MinIO object key \\
             & file\_name & VARCHAR & Original file name \\
             & file\_type & VARCHAR? & MIME type \\
             & file\_size & INT? & Size in bytes \\
  \midrule
  Submission & id & UUID (PK) & Unique identifier \\
             & post\_id & UUID (FK) & References Post (assignment) \\
             & student\_id & UUID & Student (logical FK) \\
             & content & TEXT? & Optional text content \\
             & points & INT? & Awarded grade \\
             & feedback & TEXT? & Instructor feedback \\
             & graded\_at & TIMESTAMP? & When graded \\
  \midrule
  GradeCategory & id & UUID (PK) & Unique identifier \\
                & class\_id & UUID & Classroom (logical FK) \\
                & name & VARCHAR & Category name \\
                & weight & INT & Weight percentage \\
                & position & INT & Display order \\
  \bottomrule
\end{longtable}

\begin{longtable}{@{} l l l p{5.5cm} @{}}
  \caption{Attribute dictionary --- \texttt{communication\_db} and \texttt{notifications\_db}.} \label{tab:dict-comm-notif} \\
  \toprule
  \textbf{Entity} & \textbf{Attribute} & \textbf{Type} & \textbf{Description} \\
  \midrule
  \endfirsthead
  \toprule
  \textbf{Entity} & \textbf{Attribute} & \textbf{Type} & \textbf{Description} \\
  \midrule
  \endhead
  ChatRoom & id & UUID (PK) & Unique identifier \\
           & class\_id & UUID (UNIQUE) & One chat room per classroom \\
           & enabled & BOOLEAN & Whether chat is active \\
  \midrule
  Message & id & UUID (PK) & Unique identifier \\
          & chat\_room\_id & UUID (FK) & References ChatRoom \\
          & sender\_id & UUID & Sender (logical FK) \\
          & content & TEXT? & Message text \\
          & file\_key & VARCHAR? & Attached file key \\
          & file\_name & VARCHAR? & Attached file name \\
  \midrule
  Notification & id & UUID (PK) & Unique identifier \\
               & user\_id & UUID & Target user (logical FK) \\
               & type & ENUM & Notification type (7 values) \\
               & title & VARCHAR & Notification title \\
               & body & TEXT? & Optional body text \\
               & class\_id & UUID? & Related classroom \\
               & post\_id & UUID? & Related post \\
               & is\_read & BOOLEAN & Read status \\
  \bottomrule
\end{longtable}

% ─────────────────────────────────────────────────────────────────────────────
\subsection{Interaction Modeling}
\label{sec:interaction-modeling}

The design-level sequence diagrams below detail the internal interactions within each service, showing routes, middleware, controllers, repositories, and utility classes.

\begin{figure}[H]
  \centering
  \includegraphics[width=\textwidth,height=0.85\textheight,keepaspectratio]{diagrams/seq-design-grade.png}
  \caption{Design sequence diagram --- Grading a Submission.}
  \label{fig:seq-design-grade}
\end{figure}

\begin{figure}[H]
  \centering
  \includegraphics[width=\textwidth,height=0.85\textheight,keepaspectratio]{diagrams/seq-chat.png}
  \caption{System sequence diagram --- Real-Time Chat (WebSocket).}
  \label{fig:seq-chat}
\end{figure}

\begin{figure}[H]
  \centering
  \rotatebox{90}{\includegraphics[width=0.85\textheight,height=\textwidth,keepaspectratio]{diagrams/seq-search.png}}
  \caption{System sequence diagram --- Classroom Search.}
  \label{fig:seq-search}
\end{figure}

% =============================================================================
\section*{Conclusion}

This chapter presented the system design of \eduspace{}.
The microservices architectural paradigm was justified based on domain alignment, independent data ownership, and technology suitability.
The overall architecture was described in detail, including the seven backend services, the API Gateway, the NGINX reverse proxy, the synchronous (REST) and asynchronous (RabbitMQ) communication patterns, and the infrastructure components managed by Docker Compose.
The object-oriented design artifacts---class diagrams, the persistence model, the attribute dictionary, and interaction diagrams---formalize the internal structure and dynamic behavior of the system.

The next chapter covers the implementation and deployment, detailing the technology stack, development environment, DevOps practices, security measures, and the user interfaces.
